\documentclass[11pt,letterpaper]{article}

\usepackage[utf8]{inputenc}
\usepackage[T1]{fontenc}
\usepackage[spanish]{babel}
\usepackage{geometry}
\usepackage{array}
\usepackage{longtable}
\usepackage{hyperref}
\usepackage{xcolor}
\usepackage{listings}
\usepackage{enumitem}
\usepackage{float}

\geometry{margin=2.2cm}
\setlength{\parskip}{0.45em}
\setlength{\parindent}{0pt}
\renewcommand{\arraystretch}{1.12}
\setlist[itemize]{noitemsep, topsep=2pt}
\setlist[enumerate]{noitemsep, topsep=2pt}
\sloppy

\hypersetup{
  colorlinks=true,
  linkcolor=blue,
  urlcolor=blue,
  pdftitle={Informe tecnico - Sistema de Monitoreo Remoto},
  pdfauthor={Grupo de Redes de Computadores}
}

\lstdefinestyle{terminal}{
  basicstyle=\ttfamily\footnotesize,
  breaklines=true,
  frame=single,
  columns=fullflexible,
  backgroundcolor=\color{gray!8}
}

\lstdefinestyle{json}{
  basicstyle=\ttfamily\footnotesize,
  breaklines=true,
  frame=single,
  columns=fullflexible,
  backgroundcolor=\color{gray!8}
}

\title{\textbf{Sistema de Monitoreo Remoto de Servidores Simulados}\\
\large Informe tecnico final}
\author{Curso: Redes de Computadores\\
Docente: Diego Arroyo Navarrete\\
Integrantes: \textit{[completar por el grupo]}}
\date{\textit{[completar fecha]}}

\begin{document}
\maketitle

\begin{abstract}
Este informe presenta el diseno, implementacion y validacion de un sistema de
monitoreo remoto de servidores simulados. El proyecto usa comunicacion TCP real,
clientes persistentes, autenticacion con clave precompartida por nodo, desafio
con nonce, cifrado simetrico de mensajes de aplicacion, persistencia SQLite y un
dashboard web para demostracion en tiempo real. El servidor detecta anomalias en
CPU, RAM, latencia, servicio web y eventos, emite comandos correctivos y registra
ACKs para mantener trazabilidad. El resultado es un prototipo reproducible,
defendible desde redes de computadores y acompanado por pruebas automatizadas,
escenarios de demo, evidencia SQLite, Wireshark/tshark y Nmap.
\end{abstract}

\tableofcontents
\newpage

\section{Resumen ejecutivo}

El sistema implementado resuelve un caso acotado de monitoreo remoto: varios
nodos simulados reportan periodicamente metricas operacionales a un servidor
central. El servidor recibe esas metricas por TCP, autentica al nodo, procesa el
mensaje, persiste la informacion y, cuando corresponde, responde con comandos de
mitigacion. El cliente no solo recibe comandos, sino que modifica su estado
simulado para que las metricas posteriores muestren recuperacion progresiva.

La version final incorpora una decision importante de seguridad: antes de aceptar
metricas, el servidor exige un handshake basado en PSK. El cliente se identifica
con su \texttt{node\_id}; si existe una clave asociada, el servidor envia un nonce
y el cliente responde una prueba HMAC. Despues de validar esa prueba, ambos lados
derivan una clave de sesion y los mensajes de aplicacion viajan dentro de frames
\texttt{secure} cifrados y autenticados.

El proyecto esta pensado para una defensa practica. La demostracion se puede
hacer desde el dashboard, ejecutando escenarios como CPU alta o multi-nodo. Al
mismo tiempo, la evidencia tecnica se obtiene desde SQLite, pruebas unitarias e
integracion, capturas de trafico y escaneo de puerto. Asi, el sistema no depende
solo de una interfaz visual: cada comportamiento mostrado en la demo tiene una
huella persistida o verificable.

\begin{longtable}{|p{0.28\textwidth}|p{0.62\textwidth}|}
\hline
\textbf{Aspecto} & \textbf{Estado final} \\
\hline
Comunicacion & TCP real, puerto 5000 por defecto, clientes persistentes. \\
\hline
Autenticacion & PSK por nodo, 32 claves listas: \texttt{node-01} a \texttt{node-32}. \\
\hline
Confidencialidad & Frames de aplicacion cifrados despues del handshake. \\
\hline
Servidor & Multicliente con hilos, validacion, deteccion y comandos. \\
\hline
Cliente & Reconexion automatica y mitigacion simulada local. \\
\hline
Persistencia & SQLite con tablas \texttt{metrics}, \texttt{commands} y \texttt{acks}. \\
\hline
Dashboard & API local y frontend en tiempo real para demo. \\
\hline
Validacion & 63 pruebas automatizadas, todas aprobadas en la ultima ejecucion. \\
\hline
\end{longtable}

\newpage
\section{Problema, alcance y objetivos}

\subsection{Problema abordado}

En muchos entornos operacionales se requiere observar el estado de equipos o
servicios sin acceso fisico directo. Aunque este proyecto no controla hardware
real, modela el problema central de redes: existe un conjunto de nodos remotos,
un canal de comunicacion, un servidor central y un flujo de mensajes que debe ser
validado, procesado y registrado. En ese contexto, no basta con imprimir metricas
en pantalla; el sistema debe demostrar conexion real, errores manejados,
persistencia y evidencia reproducible.

El problema se vuelve mas interesante cuando los nodos no solo reportan datos,
sino que tambien reciben ordenes. Por ejemplo, si un nodo reporta CPU alta, el
servidor debe emitir una accion correctiva y esperar una confirmacion. Ese flujo
obliga a resolver bidireccionalidad, correlacion de comandos con ACKs y estados
como \texttt{pending}, \texttt{confirmed}, \texttt{failed} o \texttt{timed\_out}.

\subsection{Alcance funcional}

El alcance del sistema incluye simulacion de nodos, comunicacion TCP, seguridad
basica mejorada con PSK, persistencia, visualizacion y pruebas. No se busca
implementar un producto de monitoreo industrial completo ni un protocolo
criptografico productivo equivalente a TLS. El foco esta en construir un
prototipo claro, ejecutable localmente y defendible desde los contenidos del
curso.

\subsection{Objetivo general}

Disenar e implementar un sistema cliente-servidor que permita monitorear nodos
simulados, registrar actividad, detectar anomalias, emitir comandos correctivos
y demostrar su funcionamiento mediante pruebas automatizadas y analisis de red.

\subsection{Objetivos especificos}

\begin{itemize}
  \item Implementar comunicacion TCP real entre clientes y servidor.
  \item Mantener clientes persistentes con reconexion automatica.
  \item Autenticar nodos mediante PSK, nonce y prueba HMAC.
  \item Cifrar los mensajes de aplicacion despues del handshake.
  \item Validar estructura, tipos y rangos de las metricas recibidas.
  \item Persistir metricas, comandos y ACKs en SQLite.
  \item Visualizar estado y eventos en un dashboard local.
  \item Entregar escenarios reproducibles y evidencia para la defensa.
\end{itemize}

\newpage
\section{Integrantes y roles}

La distribucion de roles permite explicar el proyecto como un trabajo por capas:
arquitectura, cliente, servidor/frontend y persistencia/evidencia. La tabla debe
ser completada con los nombres reales del grupo antes de entregar el PDF.

\begin{longtable}{|p{0.22\textwidth}|p{0.26\textwidth}|p{0.42\textwidth}|}
\hline
\textbf{Integrante} & \textbf{Rol principal} & \textbf{Aporte esperado o realizado} \\
\hline
\textit{[Nombre]} & Arquitectura y documentacion & Diseno general, diagrama, contrato tecnico, coherencia del informe y preparacion de defensa. \\
\hline
\textit{[Nombre]} & Cliente simulado & Cliente TCP persistente, generacion de metricas, modos anomalos, reconexion y mitigacion local. \\
\hline
\textit{[Nombre]} & Servidor y dashboard & Servidor multicliente, sesiones, reglas de anomalia, comandos y visualizacion web. \\
\hline
\textit{[Nombre]} & Persistencia, pruebas y red & SQLite, pruebas automatizadas, escenarios, Wireshark/tshark, Nmap y evidencia final. \\
\hline
\end{longtable}

\subsection{Relacion entre roles y modulos}

Aunque los roles se separan para organizar el trabajo, el proyecto final exige
integracion. El cliente no se valida sin servidor; el servidor no es demostrable
sin persistencia; el dashboard no debe inventar estado, sino leer datos reales; y
la documentacion debe coincidir con el comportamiento implementado. Por eso, la
entrega se revisa como sistema completo y no como archivos aislados.

\begin{longtable}{|p{0.28\textwidth}|p{0.58\textwidth}|}
\hline
\textbf{Modulo} & \textbf{Responsabilidad dentro del sistema} \\
\hline
\texttt{client/tcp\_client.py} & Nodo simulado persistente, metricas, comandos, ACKs y reconexion. \\
\hline
\texttt{server/client\_session.py} & Handshake seguro, procesamiento de metricas y ACKs, envio de errores y comandos. \\
\hline
\texttt{shared/secure\_channel.py} & Handshake PSK, derivacion de clave, cifrado e integridad de frames. \\
\hline
\texttt{storage/store.py} & Persistencia SQLite para metricas, comandos y confirmaciones. \\
\hline
\texttt{frontend/dashboard\_server.py} & API HTTP local, estado consolidado y ejecucion de escenarios desde la UI. \\
\hline
\texttt{tests/} & Evidencia automatizada de comportamiento esperado. \\
\hline
\end{longtable}

\newpage
\section{Arquitectura general}

La arquitectura sigue un modelo cliente-servidor centralizado. Los nodos abren
conexiones TCP hacia el servidor. El servidor acepta multiples conexiones y crea
una sesion por cliente. Cada sesion realiza autenticacion PSK antes de procesar
datos de aplicacion. Una vez autenticada la conexion, las metricas se descifran,
validan, persisten y evaluan contra reglas de anomalia. Si alguna regla se
cumple, se genera un comando, se persiste como pendiente y se envia al cliente.

\begin{figure}[H]
\centering
\fbox{\begin{minipage}{0.95\textwidth}
\centering
\begin{tabular}{ccc}
\textbf{Clientes TCP} & & \textbf{Servidor TCP} \\
\begin{tabular}{|c|}
\hline
node-01 ... node-32\\
client/tcp\_client.py\\
metricas + ACKs\\
\hline
\end{tabular}
& \(\xrightarrow{\quad TCP / PSK + frames secure \quad}\) &
\begin{tabular}{|c|}
\hline
connection\_manager.py\\
client\_session.py\\
server\_state.py\\
command\_dispatcher.py\\
\hline
\end{tabular}
\\[1.4em]
& & \(\downarrow\) \\
& &
\begin{tabular}{|c|}
\hline
SQLite\\
metrics / commands / acks\\
\hline
\end{tabular}
\\[1.4em]
& & \(\downarrow\) \\
& &
\begin{tabular}{|c|}
\hline
Dashboard HTTP\\
/api/state / /api/status / /api/scenario\\
\hline
\end{tabular}
\end{tabular}
\end{minipage}}
\caption{Arquitectura logica del sistema.}
\end{figure}

\subsection{Capas del sistema}

\textbf{Capa de comunicacion.} Usa sockets TCP. El servidor escucha en el puerto
5000 por defecto y los clientes se conectan a ese puerto. TCP entrega
confiabilidad de transporte, orden de bytes y deteccion de cierre de conexion,
pero no entrega por si solo autenticacion ni cifrado.

\textbf{Capa de seguridad de aplicacion.} Antes de enviar metricas, el cliente
ejecuta un handshake PSK. Esta capa esta implementada en
\texttt{shared/secure\_channel.py}. Su tarea es envolver el socket en una
abstraccion que envia y recibe objetos JSON cifrados.

\textbf{Capa de dominio.} Define que es una metrica valida, cuando existe una
anomalia, que comando corresponde y como se interpreta un ACK. Esta logica vive
principalmente en \texttt{server/client\_session.py},
\texttt{server/server\_state.py} y \texttt{client/tcp\_client.py}.

\textbf{Capa de persistencia y observabilidad.} SQLite registra la historia del
sistema. El dashboard consulta esa historia y la transforma en tablas, graficos y
eventos utiles para la presentacion.

\newpage
\section{Diseno del protocolo}

\subsection{Formato de linea}

El sistema usa JSON Lines: cada mensaje serializado termina en salto de linea.
Esta decision simplifica el framing sobre TCP, ya que TCP no preserva fronteras
de mensajes. Con JSON Lines, el receptor acumula bytes hasta encontrar
\texttt{\textbackslash n}; luego parsea esa linea como objeto JSON.

\subsection{Handshake PSK}

El handshake ocurre en texto JSON antes de activar el canal cifrado. El cliente
envia su identificador de nodo, el servidor busca la PSK asociada y responde con
un nonce. El cliente genera un nonce propio y calcula una prueba HMAC usando la
PSK, el \texttt{node\_id}, el nonce del servidor y el nonce del cliente.

\begin{lstlisting}[style=json]
Cliente -> Servidor:
{"type":"hello","node_id":"node-01"}

Servidor -> Cliente:
{"type":"challenge","node_id":"node-01","nonce":"..."}

Cliente -> Servidor:
{"type":"challenge_response","node_id":"node-01",
 "client_nonce":"...","proof":"..."}

Servidor -> Cliente:
{"type":"ready","node_id":"node-01"}
\end{lstlisting}

Si el nodo no existe, si falta un campo o si el HMAC no coincide, el servidor
responde con \texttt{AUTH\_FAILED} o un error equivalente y la sesion no queda
habilitada para recibir metricas.

\subsection{Derivacion de clave de sesion}

Una vez autenticada la identidad, ambas partes derivan una clave de sesion desde
la PSK y los nonces. Esto evita usar directamente la PSK como clave de cifrado de
todos los mensajes. La clave de sesion se separa internamente en una clave para
generar flujo de cifrado y otra para etiquetas de integridad.

\subsection{Frame seguro}

Despues del mensaje \texttt{ready}, el payload de aplicacion no viaja como JSON
plano. Viaja dentro de un frame \texttt{secure}.

\begin{lstlisting}[style=json]
{
  "type": "secure",
  "seq": 0,
  "nonce": "...",
  "ciphertext": "...",
  "tag": "..."
}
\end{lstlisting}

El campo \texttt{seq} evita reordenamientos o repeticiones simples dentro de la
sesion. El \texttt{nonce} se usa para generar un flujo distinto por mensaje. El
\texttt{ciphertext} contiene el JSON de aplicacion cifrado. El \texttt{tag}
autentica la secuencia, nonce y ciphertext, por lo que una alteracion del frame
produce error de protocolo.

\newpage
\section{Mensajes de aplicacion}

\subsection{Metric}

El mensaje \texttt{metric} representa el estado reportado por un nodo. Incluye
identificador, numero de secuencia y variables observables. Aunque viaja cifrado,
su estructura logica es la siguiente:

\begin{lstlisting}[style=json]
{
  "type": "metric",
  "node_id": "node-01",
  "seq": 1,
  "cpu": 45.0,
  "ram": 60.0,
  "latency_ms": 35,
  "service_web": "ok",
  "event_log": "normal",
  "token": "node-01-secret",
  "scenario": "normal",
  "anomaly_active": false,
  "mitigation_active": false,
  "mitigation_type": ""
}
\end{lstlisting}

El servidor valida que el \texttt{node\_id} del mensaje coincida con el nodo
autenticado en el canal seguro. Tambien valida tipos y rangos: CPU y RAM deben
estar entre 0 y 100, la latencia no puede ser negativa y el servicio web solo
acepta \texttt{ok} o \texttt{falla}.

\subsection{Command}

Cuando el servidor detecta una anomalia, responde con un comando correctivo.

\begin{lstlisting}[style=json]
{
  "type": "command",
  "command_id": 1,
  "action": "reduce_cpu",
  "reason": "cpu above 90"
}
\end{lstlisting}

El \texttt{command\_id} permite correlacionar la orden con el ACK posterior. El
comando se persiste con estado inicial \texttt{pending}. Cuando llega el ACK, el
estado cambia a \texttt{confirmed} o \texttt{failed}.

\subsection{ACK}

El ACK confirma si el cliente aplico o no la orden. En la demo, el cliente agrega
informacion de estado antes y despues de aplicar la mitigacion.

\begin{lstlisting}[style=json]
{
  "type": "ack",
  "node_id": "node-01",
  "command_id": 1,
  "status": "applied",
  "command": "reduce_cpu",
  "before": {"cpu": null, "mitigation_active": false},
  "after": {"cpu": 95.0, "mitigation_active": true}
}
\end{lstlisting}

\subsection{Error}

El mensaje \texttt{error} se usa para indicar fallas de autenticacion,
validacion o protocolo. En una implementacion real, convendria separar errores
expuestos al cliente de errores internos detallados, pero para la demo resulta
util mostrar el motivo basico del rechazo.

\newpage
\section{Autenticacion, claves y limitaciones criptograficas}

\subsection{Claves precompartidas}

La demo deja configuradas 32 claves, desde \texttt{node-01} hasta
\texttt{node-32}. El patron usado es \texttt{node-XX-secret}. Por ejemplo:

\begin{longtable}{|p{0.26\textwidth}|p{0.48\textwidth}|}
\hline
\textbf{Nodo} & \textbf{PSK de demo} \\
\hline
\texttt{node-01} & \texttt{node-01-secret} \\
\hline
\texttt{node-02} & \texttt{node-02-secret} \\
\hline
\texttt{node-03} & \texttt{node-03-secret} \\
\hline
\texttt{...} & \texttt{...} \\
\hline
\texttt{node-32} & \texttt{node-32-secret} \\
\hline
\end{longtable}

Estas claves se encuentran en el codigo para facilitar ejecucion y evaluacion.
En un despliegue real, las claves no deberian estar escritas directamente en el
repositorio; deberian provenir de variables de entorno, un gestor de secretos o
un archivo protegido fuera del control de versiones.

\subsection{Por que usar nonce}

El nonce evita que la prueba de identidad sea siempre igual. Si el cliente
enviara simplemente la clave o un hash fijo, una captura de red podria reutilizar
ese valor. Con el nonce, el servidor exige una prueba distinta en cada conexion.
El cliente demuestra que conoce la PSK sin enviarla directamente.

\subsection{Alcance del cifrado}

La implementacion usa biblioteca estandar de Python. Se construye un flujo de
cifrado derivado de HMAC-SHA256 y se agrega una etiqueta HMAC de integridad. Esto
es suficiente para una demo academica reproducible sin dependencias externas, y
permite observar en Wireshark que los payloads de aplicacion ya no aparecen en
claro despues del handshake.

Sin embargo, no se presenta como reemplazo productivo de TLS ni de una AEAD
auditada como AES-GCM o ChaCha20-Poly1305. La conclusion correcta para la defensa
es que el proyecto demuestra conceptos de autenticacion, derivacion de clave,
confidencialidad e integridad, pero que un sistema real deberia usar TLS o una
biblioteca criptografica especializada.

\subsection{Controles implementados}

\begin{itemize}
  \item Nodo desconocido: se rechaza en handshake.
  \item Prueba HMAC invalida: se rechaza antes de crear canal seguro.
  \item Mensaje fuera de secuencia: se rechaza el frame.
  \item Tag invalido: se rechaza el frame por integridad.
  \item \texttt{node\_id} distinto dentro del canal: se rechaza el mensaje.
\end{itemize}

\newpage
\section{Servidor TCP}

\subsection{Aceptacion de conexiones}

El servidor escucha conexiones entrantes en el host y puerto configurados. Por
defecto se usa el puerto 5000. Cada conexion aceptada se atiende en una sesion
independiente, lo que permite que varios clientes reporten metricas al mismo
tiempo sin bloquearse entre si.

\subsection{Sesion de cliente}

La clase de sesion tiene tres responsabilidades principales. Primero, establece
el canal seguro mediante \texttt{server\_handshake}. Segundo, procesa mensajes de
aplicacion descifrados. Tercero, envia respuestas cifradas, como comandos o
errores.

El flujo interno simplificado es:

\begin{enumerate}
  \item aceptar socket ya conectado;
  \item leer \texttt{hello};
  \item autenticar PSK con nonce;
  \item crear \texttt{SecureSocket};
  \item leer frames seguros en bucle;
  \item despachar \texttt{metric} o \texttt{ack};
  \item cerrar y marcar desconexion al terminar.
\end{enumerate}

\subsection{Validacion de metricas}

La validacion no se limita a parsear JSON. El servidor revisa que las metricas
tengan valores razonables. Por ejemplo, CPU=150 se rechaza aunque sea JSON valido.
Esto evita que datos corruptos entren a la base y que el dashboard muestre estados
imposibles.

\begin{longtable}{|p{0.24\textwidth}|p{0.58\textwidth}|}
\hline
\textbf{Campo} & \textbf{Regla de validacion} \\
\hline
\texttt{node\_id} & Debe existir y coincidir con el nodo autenticado. \\
\hline
\texttt{seq} & Entero mayor o igual a 0. \\
\hline
\texttt{cpu} & Numero entre 0 y 100. \\
\hline
\texttt{ram} & Numero entre 0 y 100. \\
\hline
\texttt{latency\_ms} & Numero mayor o igual a 0. \\
\hline
\texttt{service\_web} & Valor \texttt{ok} o \texttt{falla}. \\
\hline
\end{longtable}

\subsection{Cooldown anti-spam}

El servidor evita emitir comandos duplicados de manera excesiva. Si un nodo ya
tiene una accion pendiente o confirmada recientemente, el sistema suprime el
comando repetido. Esto mejora la demo: no se llena el historial con decenas de
\texttt{reduce\_cpu} para la misma anomalia mientras el cliente aun se recupera.

\newpage
\section{Cliente TCP y simulacion de nodos}

\subsection{Cliente persistente}

El cliente real se ejecuta con \texttt{python -m client.tcp\_client}. Mantiene una
conexion TCP abierta, envia metricas cada cierto intervalo y escucha comandos del
servidor. Si la conexion se pierde, espera cinco segundos e intenta reconectar,
conservando el contador de secuencia.

\begin{lstlisting}[style=terminal]
python3 -m client.tcp_client --node-id node-01 --mode normal
python3 -m client.tcp_client --node-id node-02 --mode high-cpu
python3 -m client.tcp_client --node-id node-03 --mode high-latency --interval 3.0
\end{lstlisting}

\subsection{Modos de simulacion}

\begin{longtable}{|p{0.25\textwidth}|p{0.62\textwidth}|}
\hline
\textbf{Modo} & \textbf{Comportamiento} \\
\hline
\texttt{normal} & Reporta valores base: CPU 35, RAM 45, latencia 40 ms. \\
\hline
\texttt{high-cpu} & Reporta CPU 95 y dispara \texttt{reduce\_cpu}. \\
\hline
\texttt{high-ram} & Reporta RAM 94 y dispara \texttt{reduce\_ram}. \\
\hline
\texttt{high-latency} & Reporta latencia 350 ms y dispara \texttt{fix\_latency}. \\
\hline
\texttt{service-failure} & Reporta servicio web en falla y dispara \texttt{restart\_service}. \\
\hline
\texttt{failed-event} & Reporta evento con texto \texttt{fallido} y dispara \texttt{normalize\_node}. \\
\hline
\texttt{chaos} & Cicla deterministamente por todas las anomalias. \\
\hline
\end{longtable}

\subsection{Mitigacion local}

La mitigacion es simulada, pero no es solo un ACK vacio. Cuando el cliente recibe
un comando, actualiza su estado interno. Por ejemplo, ante \texttt{reduce\_cpu},
las metricas posteriores comienzan a bajar gradualmente desde el valor anomalico
hacia el valor base. Esto permite mostrar en el dashboard una historia completa:
anomalia, comando, ACK y recuperacion.

\subsection{Reconexion}

La reconexion es importante porque demuestra comportamiento realista ante fallas.
Si se detiene el servidor, el cliente detecta error de socket, informa la perdida
de conexion y reintenta. Al volver a levantar el servidor, el cliente realiza de
nuevo el handshake PSK y continua enviando metricas.

\newpage
\section{Reglas de anomalia y comandos}

El servidor aplica reglas simples por umbral o evento. La intencion no es crear
un motor de monitoreo complejo, sino representar claramente la transformacion de
un dato recibido en una decision operacional.

\begin{longtable}{|p{0.28\textwidth}|p{0.24\textwidth}|p{0.34\textwidth}|}
\hline
\textbf{Condicion} & \textbf{Comando} & \textbf{Interpretacion} \\
\hline
\texttt{cpu > 90} & \texttt{reduce\_cpu} & El nodo reporta uso excesivo de CPU. \\
\hline
\texttt{ram > 90} & \texttt{reduce\_ram} & El nodo reporta presion de memoria. \\
\hline
\texttt{latency\_ms > 200} & \texttt{fix\_latency} & El nodo experimenta latencia elevada. \\
\hline
\texttt{service\_web == falla} & \texttt{restart\_service} & El servicio web simulado no esta operativo. \\
\hline
\texttt{event\_log} contiene \texttt{fallido} & \texttt{normalize\_node} & El log contiene una falla operacional. \\
\hline
\end{longtable}

\subsection{Estados de comando}

Cada comando emitido se registra con estado. Esto permite auditar si el servidor
solo mando ordenes o si realmente recibio confirmacion.

\begin{longtable}{|p{0.24\textwidth}|p{0.62\textwidth}|}
\hline
\textbf{Estado} & \textbf{Significado} \\
\hline
\texttt{pending} & El comando fue enviado, pero aun no existe ACK valido. \\
\hline
\texttt{confirmed} & El cliente respondio con ACK \texttt{applied}. \\
\hline
\texttt{failed} & El cliente respondio con ACK \texttt{failed}. \\
\hline
\texttt{timed\_out} & El comando expiro sin confirmacion dentro del plazo. \\
\hline
\end{longtable}

\subsection{Ejemplo narrativo}

En el escenario CPU alta, el cliente \texttt{node-02} reporta CPU=95. El servidor
valida la metrica, la persiste y detecta que supera el umbral. Luego construye un
comando \texttt{reduce\_cpu}, lo persiste como \texttt{pending} y lo envia cifrado.
El cliente recibe el comando, ajusta su estado, responde un ACK \texttt{applied} y
las metricas siguientes muestran CPU bajando progresivamente.

\newpage
\section{Persistencia SQLite}

La persistencia convierte la demo en evidencia. Sin base de datos, el sistema
podria mostrar eventos momentaneos en consola, pero seria dificil auditar que
ocurrio. SQLite permite consultar metricas, comandos y ACKs despues de ejecutar
un escenario.

\subsection{Tabla metrics}

\begin{longtable}{|p{0.22\textwidth}|p{0.16\textwidth}|p{0.48\textwidth}|}
\hline
\textbf{Campo} & \textbf{Tipo} & \textbf{Descripcion} \\
\hline
\texttt{id} & INTEGER & Identificador autoincremental. \\
\hline
\texttt{node\_id} & TEXT & Nodo que reporto la metrica. \\
\hline
\texttt{seq} & INTEGER & Numero de secuencia del cliente. \\
\hline
\texttt{cpu} & REAL & Porcentaje de CPU. \\
\hline
\texttt{ram} & REAL & Porcentaje de RAM. \\
\hline
\texttt{latency\_ms} & REAL & Latencia en milisegundos. \\
\hline
\texttt{service\_web} & TEXT & Estado del servicio web. \\
\hline
\texttt{event\_log} & TEXT & Evento textual reportado. \\
\hline
\texttt{scenario} & TEXT & Modo de simulacion usado por el cliente. \\
\hline
\texttt{mitigation\_active} & INTEGER & Indica si el cliente esta mitigando. \\
\hline
\end{longtable}

\subsection{Tabla commands}

La tabla \texttt{commands} registra que orden emitio el servidor, hacia que nodo
y en que estado quedo. Su valor no es solo historico; tambien sirve para verificar
que los ACKs actualizan correctamente el estado.

\subsection{Tabla acks}

La tabla \texttt{acks} registra la confirmacion enviada por el cliente. En la
defensa, esta tabla permite mostrar que la respuesta al comando no fue solo un
evento visual del dashboard, sino una fila persistida y correlacionable.

\subsection{Consultas utiles}

\begin{lstlisting}[style=terminal]
sqlite3 data/monitor.db "SELECT node_id, seq, cpu, ram FROM metrics ORDER BY id DESC LIMIT 10;"
sqlite3 data/monitor.db "SELECT command_id, node_id, action, status FROM commands;"
sqlite3 data/monitor.db "SELECT command_id, node_id, status FROM acks;"
\end{lstlisting}

\newpage
\section{Dashboard y API HTTP}

El dashboard no reemplaza las pruebas ni la evidencia, pero hace visible el
comportamiento del sistema durante la presentacion. Sirve HTML, CSS y JavaScript
desde un servidor HTTP local y consulta endpoints de API para obtener estado.

\subsection{Fuente unica de verdad}

El endpoint principal es \texttt{/api/state}. Este endpoint construye un snapshot
desde SQLite y estado del servidor. El frontend no deberia inventar nodos activos
ni comandos; debe reflejar lo que el backend conoce.

\begin{longtable}{|p{0.28\textwidth}|p{0.58\textwidth}|}
\hline
\textbf{Endpoint} & \textbf{Uso} \\
\hline
\texttt{/} & Sirve la interfaz HTML. \\
\hline
\texttt{/app.js} y \texttt{/styles.css} & Archivos estaticos del frontend. \\
\hline
\texttt{/api/status} & Indica si el servidor TCP esta alcanzable. \\
\hline
\texttt{/api/state} & Devuelve snapshot completo del sistema. \\
\hline
\texttt{/api/metrics} & Expone metricas persistidas. \\
\hline
\texttt{/api/events} & Expone eventos combinados de comandos y ACKs. \\
\hline
\texttt{/api/scenario} & Ejecuta clientes reales para escenarios de demo. \\
\hline
\end{longtable}

\subsection{Escenarios desde la interfaz}

Los botones del dashboard ejecutan clientes reales contra el servidor TCP activo.
Esto es relevante: la demo no simula datos dentro del navegador, sino que dispara
el flujo real de red, autenticacion, cifrado, persistencia y lectura posterior.

\subsection{Lectura visual durante la defensa}

Durante la presentacion se recomienda mostrar:

\begin{itemize}
  \item indicador de servidor TCP online;
  \item nodo activo y ultima metrica;
  \item grafico de CPU/RAM/latencia;
  \item evento de comando emitido;
  \item ACK recibido;
  \item recuperacion de la metrica despues de la mitigacion.
\end{itemize}

\newpage
\section{Escenarios de demostracion}

\subsection{Escenario CPU alta}

Este es el escenario principal porque permite observar todo el flujo en poco
tiempo. Se ejecuta \texttt{high-cpu}, el cliente reporta CPU=95, el servidor envia
\texttt{reduce\_cpu}, el cliente confirma y luego la CPU baja gradualmente.

\begin{lstlisting}[style=terminal]
python3 -m client.tcp_client --node-id node-02 --mode high-cpu --interval 3.0
\end{lstlisting}

Resultado esperado:

\begin{enumerate}
  \item aparece el nodo en el dashboard;
  \item la serie de CPU sube;
  \item se registra comando \texttt{reduce\_cpu};
  \item aparece ACK \texttt{applied};
  \item la CPU vuelve hacia rango normal.
\end{enumerate}

\subsection{Escenario multi-nodo}

El escenario multi-nodo levanta varios clientes con modos distintos. Su objetivo
es demostrar concurrencia y diferenciacion por nodo.

\begin{longtable}{|p{0.22\textwidth}|p{0.28\textwidth}|p{0.34\textwidth}|}
\hline
\textbf{Nodo} & \textbf{Modo} & \textbf{Efecto} \\
\hline
\texttt{node-01} & \texttt{normal} & Metricas base. \\
\hline
\texttt{node-02} & \texttt{high-cpu} & CPU alta. \\
\hline
\texttt{node-03} & \texttt{high-latency} & Latencia alta. \\
\hline
\texttt{node-04} & \texttt{high-ram} & RAM alta. \\
\hline
\texttt{node-05} & \texttt{service-failure} & Servicio en falla. \\
\hline
\texttt{node-06} & \texttt{failed-event} & Evento fallido. \\
\hline
\texttt{node-07} & \texttt{chaos} & Ciclo de anomalias. \\
\hline
\end{longtable}

\subsection{Escenario nodo sin PSK valida}

Para demostrar autenticacion, se puede intentar conectar un nodo no configurado:

\begin{lstlisting}[style=terminal]
python3 -m client.tcp_client --node-id node-99 --mode normal
\end{lstlisting}

El resultado esperado es rechazo en el handshake. No se deben persistir metricas
porque el canal seguro nunca queda establecido.

\newpage
\section{Pruebas automatizadas}

La suite se ejecuta con:

\begin{lstlisting}[style=terminal]
python3 -m unittest discover -s tests -v
\end{lstlisting}

Estado actual verificado: \textbf{63 pruebas automatizadas, todas pasan}.

\subsection{Cobertura por area}

\begin{longtable}{|p{0.25\textwidth}|p{0.55\textwidth}|p{0.12\textwidth}|}
\hline
\textbf{Area} & \textbf{Que se verifica} & \textbf{Estado} \\
\hline
Cliente & Construccion de metricas, modos anomalos, mitigacion y serializacion. & OK \\
\hline
Canal seguro & Handshake PSK, nodo desconocido y \texttt{node-32}. & OK \\
\hline
Servidor & Estado de comandos, cooldown, expiracion y confirmaciones. & OK \\
\hline
Persistencia & Insercion de metricas, comandos, ACKs y cambio de estado. & OK \\
\hline
Dashboard & Endpoints, archivos estaticos, escenarios y estado consolidado. & OK \\
\hline
Higiene repo & Ausencia de dependencias o rutas removidas. & OK \\
\hline
\end{longtable}

\subsection{Por que importan estas pruebas}

Las pruebas no son solo una formalidad. En este proyecto cubren errores que
podrian afectar directamente la demo: comandos que quedan eternamente
\texttt{pending}, dashboard que marca fracaso aunque el cliente genero datos,
cliente que no se recupera despues de una mitigacion o servidor que acepta datos
sin autenticacion correcta. La suite reduce el riesgo de que un cambio de ultima
hora rompa el flujo presentado.

\subsection{Prueba relevante agregada para 32 claves}

Se agrego una prueba que autentica \texttt{node-32}. Esta prueba comprueba que el
ultimo nodo del rango configurado puede completar el handshake PSK, no solo que
existe una cadena escrita en documentacion.

\newpage
\section{Evidencia de red: Wireshark y Nmap}

\subsection{Wireshark o tshark}

La captura debe hacerse filtrando el puerto TCP 5000. Se espera observar primero
el handshake TCP, luego el handshake PSK y despues frames \texttt{secure}.

\begin{lstlisting}[style=terminal]
tcp.port == 5000
\end{lstlisting}

Durante el handshake PSK se pueden ver mensajes JSON como \texttt{hello},
\texttt{challenge}, \texttt{challenge\_response} y \texttt{ready}. Despues del
\texttt{ready}, los mensajes reales de aplicacion no deben verse como
\texttt{metric}, \texttt{command} o \texttt{ack} en claro; deben aparecer como
campos \texttt{ciphertext} dentro de frames \texttt{secure}.

\begin{longtable}{|p{0.18\textwidth}|p{0.68\textwidth}|}
\hline
\textbf{Fase} & \textbf{Evidencia esperada} \\
\hline
TCP & Paquetes \texttt{SYN}, \texttt{SYN-ACK}, \texttt{ACK}. \\
\hline
PSK & Mensajes \texttt{hello}, \texttt{challenge}, \texttt{challenge\_response}, \texttt{ready}. \\
\hline
Canal seguro & Frames con \texttt{seq}, \texttt{nonce}, \texttt{ciphertext}, \texttt{tag}. \\
\hline
Reconexion & Nuevo handshake TCP y nuevo handshake PSK. \\
\hline
\end{longtable}

\subsection{Nmap}

Nmap permite demostrar que el puerto del servidor esta abierto cuando el servicio
esta ejecutandose.

\begin{lstlisting}[style=terminal]
nmap -p 5000 127.0.0.1
\end{lstlisting}

Resultado esperado:

\begin{lstlisting}[style=terminal]
5000/tcp open
\end{lstlisting}

\newpage
\section{Ejecucion reproducible}

\subsection{Servidor}

\begin{lstlisting}[style=terminal]
python3 -m server.tcp_server
\end{lstlisting}

El servidor escucha en \texttt{0.0.0.0:5000} por defecto. Tambien puede cambiarse
el host o puerto mediante configuracion.

\subsection{Dashboard}

\begin{lstlisting}[style=terminal]
python3 frontend/dashboard_server.py
\end{lstlisting}

Luego se abre:

\begin{lstlisting}[style=terminal]
http://127.0.0.1:8080
\end{lstlisting}

\subsection{Proyecto completo}

\begin{lstlisting}[style=terminal]
bash scripts/run_project.sh
\end{lstlisting}

Este comando levanta servidor y dashboard. Para una demo guiada tambien puede
usarse:

\begin{lstlisting}[style=terminal]
make present
\end{lstlisting}

\subsection{Prueba manual recomendada}

\begin{enumerate}
  \item Levantar servidor y dashboard.
  \item Abrir \texttt{http://127.0.0.1:8080}.
  \item Ejecutar escenario CPU alta.
  \item Confirmar comando \texttt{reduce\_cpu} y ACK.
  \item Ver que CPU baja progresivamente.
  \item Ejecutar escenario multi-nodo.
  \item Revisar SQLite y captura Wireshark.
\end{enumerate}

\newpage
\section{Analisis de confiabilidad}

\subsection{Tolerancia a desconexion}

El cliente esta disenado para reconectarse. Esto significa que una caida temporal
del servidor no obliga a reiniciar manualmente todos los clientes. En la defensa,
se puede detener el servidor, observar el mensaje de reconexion del cliente y
volver a levantar el servicio.

\subsection{Manejo de errores}

El servidor maneja errores comunes: JSON invalido durante handshake, nodo
desconocido, respuesta HMAC incorrecta, frames fuera de secuencia, mensajes con
campos faltantes y metricas fuera de rango. La idea es que un cliente defectuoso
no bote el proceso principal.

\subsection{Persistencia como recuperacion parcial}

SQLite no convierte el sistema en altamente disponible, pero conserva evidencia.
Si la interfaz se recarga o se reinicia el dashboard, el historial persistido
permite reconstruir metricas recientes, comandos y ACKs.

\subsection{Riesgos conocidos}

\begin{itemize}
  \item El servidor sigue siendo un punto central unico.
  \item Las claves de demo estan en el repositorio.
  \item El cifrado implementado es academico y no sustituye TLS.
  \item La ejecucion intensiva de muchos clientes depende de recursos locales.
  \item Capturas en loopback pueden variar segun sistema operativo.
\end{itemize}

\newpage
\section{Limitaciones y trabajo futuro}

\subsection{Limitaciones}

El proyecto cumple su objetivo academico, pero tiene limites claros. El mas
importante es que el esquema criptografico se implementa con biblioteca estandar
para evitar dependencias externas. Esto ayuda a la reproducibilidad, pero no debe
presentarse como una solucion productiva. Tambien existe un solo servidor central
y las PSK se generan de forma estatica.

\subsection{Trabajo futuro}

\begin{itemize}
  \item Reemplazar el cifrado de aplicacion por TLS.
  \item Cargar PSK desde variables de entorno o gestor de secretos.
  \item Agregar rotacion de claves por nodo.
  \item Incorporar autenticacion mutua con certificados.
  \item Agregar alertas por correo, webhook o mensajeria.
  \item Permitir multiples servidores o replicacion de estado.
  \item Exportar metricas a Prometheus u otro sistema de observabilidad.
  \item Mejorar dashboard con filtros historicos y descarga de evidencia.
\end{itemize}

\subsection{Decision consciente de alcance}

No se implementaron todas estas mejoras porque el objetivo de la entrega es
mostrar dominio de redes, sockets, protocolo, persistencia y evidencia. Agregar
demasiadas piezas externas podria hacer el proyecto menos reproducible para la
evaluacion. La solucion final mantiene un equilibrio entre funcionalidad real y
claridad pedagogica.

\newpage
\section{Guion de defensa}

La presentacion debe combinar explicacion y demo. No conviene separar por completo
la demo del relato tecnico; es mejor explicar cada componente justo cuando se ve
su efecto.

\begin{enumerate}
  \item Presentar el problema: monitoreo remoto y reaccion ante anomalias.
  \item Mostrar arquitectura: clientes, servidor, SQLite y dashboard.
  \item Explicar protocolo: TCP, PSK, nonce y frames cifrados.
  \item Levantar dashboard y mostrar servidor online.
  \item Ejecutar CPU alta y narrar metrica, comando, ACK y recuperacion.
  \item Mostrar SQLite o eventos persistidos.
  \item Ejecutar multi-nodo si hay tiempo.
  \item Mostrar resultado de pruebas automatizadas.
  \item Mostrar evidencia Wireshark/Nmap.
  \item Cerrar con limitaciones y mejoras futuras.
\end{enumerate}

\subsection{Frases tecnicas utiles}

\begin{itemize}
  \item ``TCP nos da transporte confiable, pero la autenticacion y cifrado se implementan en capa de aplicacion''.
  \item ``El servidor no acepta metricas antes de validar la PSK del nodo''.
  \item ``El dashboard no inventa datos: lee estado construido desde SQLite y el servidor''.
  \item ``El ACK actualiza el estado del comando, por eso no queda eternamente pending''.
  \item ``La evidencia no es solo visual; tambien queda en pruebas, base de datos y captura de red''.
\end{itemize}

\newpage
\section{Entregables}

El repositorio contiene los elementos necesarios para reproducir y defender el
proyecto.

\begin{longtable}{|p{0.3\textwidth}|p{0.56\textwidth}|}
\hline
\textbf{Entregable} & \textbf{Ubicacion o descripcion} \\
\hline
Codigo fuente & Directorios \texttt{client/}, \texttt{server/}, \texttt{shared/}, \texttt{storage/}, \texttt{frontend/}. \\
\hline
Contrato tecnico & \texttt{docs/contract\_v1.md}. \\
\hline
Arquitectura & \texttt{docs/architecture.md}. \\
\hline
Escenarios & \texttt{docs/scenarios.md}. \\
\hline
Guia de demo & \texttt{docs/demo-flow.md}. \\
\hline
Evidencia de red & \texttt{docs/evidence/wireshark.md} y \texttt{docs/evidence/nmap.md}. \\
\hline
Pruebas & \texttt{tests/} y comando \texttt{python -m unittest discover -s tests -v}. \\
\hline
Informe final & \texttt{docs/informe-final.tex}. \\
\hline
\end{longtable}

\subsection{Antes de entregar}

\begin{itemize}
  \item Completar nombres de integrantes y fecha.
  \item Insertar capturas reales si el docente las exige dentro del PDF.
  \item Generar evidencia Wireshark/Nmap en el entorno autorizado.
  \item Ejecutar nuevamente la suite de pruebas.
  \item Confirmar que el dashboard funciona en \texttt{127.0.0.1:8080}.
\end{itemize}

\newpage
\section{Conclusiones}

El proyecto implementa un sistema completo de monitoreo remoto simulado con
comunicacion TCP real. La solucion incluye clientes persistentes, servidor
concurrente, autenticacion PSK con nonce, canal de aplicacion cifrado,
persistencia SQLite, dashboard en tiempo real y pruebas automatizadas. El flujo
de anomalia, comando, ACK y recuperacion queda visible en la demo y tambien
registrado en la base de datos.

Desde el punto de vista de redes, el proyecto permite demostrar varios conceptos:
conexion TCP, framing de mensajes sobre un stream, manejo de desconexiones,
protocolos de aplicacion, autenticacion, cifrado de payload, observacion con
Wireshark y verificacion de puerto con Nmap. Desde el punto de vista de
ingenieria, muestra separacion de responsabilidades, pruebas de integracion,
persistencia y documentacion coherente con el comportamiento real.

La principal limitacion es que el cifrado implementado es academico y no reemplaza
TLS. Aun asi, cumple el objetivo de mostrar como un cliente prueba su identidad
sin enviar la clave directamente y como los mensajes dejan de viajar en claro
despues del handshake. Para una version productiva, el siguiente paso natural es
usar TLS, cargar secretos fuera del repositorio y fortalecer la gestion de nodos.

\newpage
\section{Referencias}

\begin{itemize}
  \item RFC 793: Transmission Control Protocol.
  \item Documentacion oficial de Python: \texttt{socket}, \texttt{threading}, \texttt{sqlite3}, \texttt{http.server}, \texttt{hmac}, \texttt{hashlib}.
  \item Documentacion del proyecto: \texttt{README.md}, \texttt{docs/contract\_v1.md}, \texttt{docs/architecture.md}, \texttt{docs/scenarios.md}.
  \item Guia del Proyecto Final de Redes de Computadores 2026.
\end{itemize}

\appendix
\newpage
\section{Anexo A: comandos principales}

\begin{lstlisting}[style=terminal]
# Ejecutar pruebas
python3 -m unittest discover -s tests -v

# Levantar servidor
python3 -m server.tcp_server

# Levantar dashboard
python3 frontend/dashboard_server.py

# Cliente normal
python3 -m client.tcp_client --node-id node-01 --mode normal

# Cliente con CPU alta
python3 -m client.tcp_client --node-id node-02 --mode high-cpu

# Escaneo Nmap
nmap -p 5000 127.0.0.1
\end{lstlisting}

\section{Anexo B: checklist de evaluacion}

\begin{longtable}{|p{0.5\textwidth}|p{0.18\textwidth}|p{0.2\textwidth}|}
\hline
\textbf{Criterio} & \textbf{Estado} & \textbf{Evidencia} \\
\hline
Comunicacion TCP real & Cumplido & Servidor y clientes por socket. \\
\hline
Nodos simulados & Cumplido & Modos \texttt{normal}, \texttt{high-cpu}, etc. \\
\hline
Autenticacion & Cumplido & PSK + nonce + HMAC. \\
\hline
Cifrado & Cumplido & Frames \texttt{secure}. \\
\hline
Persistencia & Cumplido & SQLite. \\
\hline
Dashboard & Cumplido & Frontend local. \\
\hline
Pruebas & Cumplido & 63 tests. \\
\hline
Evidencia de red & Documentado & Wireshark/tshark y Nmap. \\
\hline
\end{longtable}

\end{document}
